\documentclass[12pt,a4paper]{amsart}
\usepackage{amsmath,amssymb,amsthm}
\usepackage{mathtools}
\usepackage[utf8]{inputenc}
\usepackage[T1]{fontenc}
\usepackage{hyperref}
\usepackage{enumitem,booktabs,array}
\usepackage{geometry}
\geometry{margin=2.5cm}

% -------------------------------------------------------
% Theoremumgebungen (deutsche Bezeichnungen)
% -------------------------------------------------------
\newtheorem{theorem}{Satz}[section]
\newtheorem{lemma}[theorem]{Lemma}
\newtheorem{korollar}[theorem]{Korollar}
\newtheorem{proposition}[theorem]{Proposition}
\theoremstyle{definition}
\newtheorem{definition}[theorem]{Definition}
\newtheorem{beispiel}[theorem]{Beispiel}
\newtheorem{bemerkung}[theorem]{Bemerkung}
\newtheorem{vermutung}[theorem]{Vermutung}

% -------------------------------------------------------
% Eigene Befehle
% -------------------------------------------------------
\newcommand{\N}{\mathbb{N}}
\newcommand{\Z}{\mathbb{Z}}
\newcommand{\Q}{\mathbb{Q}}
\newcommand{\R}{\mathbb{R}}
\newcommand{\C}{\mathbb{C}}
\newcommand{\F}{\mathbb{F}}
\newcommand{\lf}[1]{!{#1}}
\DeclareMathOperator{\lcm}{kgV}
\DeclareMathOperator{\ord}{ord}
\DeclareMathOperator{\ggT}{ggT}

% -------------------------------------------------------
\begin{document}
% -------------------------------------------------------

\title{Die Kurepa-Vermutung:\\
       Links-Faktoriellen, Wilsons Satz\\
       und numerische Verifikation}

\author{Michael Fuhrmann}
\address{Specialist Mathematics Project, Build 133}
\email{specialist-maths@localhost}
\date{\today}

\begin{abstract}
Die Kurepa-Vermutung, 1971 vom jugoslawischen Mathematiker \DJ{}uro Kurepa
aufgestellt, besagt, dass für jede ungerade Primzahl $p$ die Links-Faktori\-elle
$!p = \sum_{k=0}^{p-1} k!$ nicht durch $p$ teilbar ist.
Trotz ihrer elementaren Formulierung ist diese Vermutung seit mehr als
fünfzig Jahren offen.
Wir geben eine in sich geschlossene Darstellung des Problems: Definition und
elementare Eigenschaften der Links-Faktoriellen, die entscheidende Verbindung
zu Wilsons Satz, probabilistische Heuristiken, bekannte Äquivalenzen sowie
eine Übersicht zur numerischen Evidenz.
Das Specialist Mathematics Project (Build 131) hat die Vermutung für alle
$664{.}578$ Primzahlen bis $10^7$ in einer Rechnung von ca.\ $8{.}396$
Sekunden verifiziert und dabei null Verletzungen gefunden.
\end{abstract}

\maketitle
\tableofcontents

% ---------------------------------------------------------------
\section{Einleitung}
% ---------------------------------------------------------------

\DJ{}uro Kurepa (1907--1993) war ein bedeutender jugoslawischer Mathematiker,
dessen Forschungsinteressen Mengenlehre, Zahlentheorie und Kombinatorik
umfassten.
Im Jahr 1971 führte er die sogenannte \emph{Links-Faktori\-elle} $!n$ ein
und stellte sogleich eine Vermutung über ihre Reste modulo Primzahlen
auf~\cite{Kurepa1971}.
Die Aussage ist verblüffend einfach: Die Links-Faktori\-elle jeder ungeraden
Primzahl $p$ ist nicht durch $p$ teilbar.

Obwohl die Formulierung nur Schulkenntnisse der modularen Arithmetik erfordert,
hat die Kurepa-Vermutung allen Beweisversuchen über fünf Jahrzehnte widerstanden.
Sie nimmt eine eigentümliche Stellung in der Zahlentheorie ein:
Probabilistische Argumente lassen erwarten, dass Gegenbeispiele existieren
(die erwartete Anzahl verletzender Primzahlen ist unendlich, analog zur
Divergenz von $\sum 1/p$), doch die numerische Suche unter den ersten
$664{.}578$ Primzahlen hat kein einziges Gegenbeispiel gefunden.
Diese Spannung zwischen heuristischer Erwartung und rechnerischer Realität
macht das Problem besonders faszinierend.

\subsection*{Aufbau der Arbeit}

Abschnitt~\ref{sec:linksfak} führt die Links-Faktori\-elle ein und leitet
elementare Eigenschaften her.
Abschnitt~\ref{sec:vermutung} formuliert die Kurepa-Vermutung präzise und
bespricht den Ausnahmefall $p = 2$.
Abschnitt~\ref{sec:wilson} untersucht die Verbindung zu Wilsons Satz und
leitet eine äquivalente Formulierung her.
Abschnitt~\ref{sec:heuristik} präsentiert probabilistische Heuristiken.
Abschnitt~\ref{sec:aequivalenzen} sammelt weitere Äquivalenzen und
Umformulierungen.
Abschnitt~\ref{sec:verifikation} beschreibt die numerische Verifikation aus
Build 131.
Abschnitt~\ref{sec:verwandt} behandelt verwandte Probleme.
Abschnitt~\ref{sec:schluss} schliesst mit offenen Fragen.

% ---------------------------------------------------------------
\section{Die Links-Faktori\-elle}
\label{sec:linksfak}
% ---------------------------------------------------------------

\subsection{Definition}

\begin{definition}[Links-Faktori\-elle]
Für $n \in \N_0$ heißt
\[
  !n \;=\; \sum_{k=0}^{n-1} k! \;=\; 0! + 1! + 2! + \cdots + (n-1)!
\]
die \emph{Links-Faktori\-elle} von $n$.
Per Konvention gilt $!0 = 0$ (leere Summe).
\end{definition}

Die Schreibweise $!n$ mit dem Ausrufezeichen \emph{links} von $n$ wurde von
Kurepa eingeführt, in Analogie zur gewöhnlichen (rechten) Faktoriellen $n!$.
In der Literatur wird $!n$ auch als \emph{Kurepa-Funktion} bezeichnet.

\begin{bemerkung}
Die Links-Faktori\-elle ist \emph{nicht} zu verwechseln mit der
\emph{Subfaktoriellen} (Anzahl der Derangements) $D_n$,
die ebenfalls manchmal $!n$ geschrieben wird.
In dieser Arbeit bezeichnet $!n$ stets Kurepas Links-Faktori\-elle.
\end{bemerkung}

\subsection{Anfangswerte}

Die folgende Tabelle zeigt $!n$ für kleine Werte von $n$:

\begin{center}
\begin{tabular}{@{}rr@{}}
\toprule
$n$ & $!n$ \\
\midrule
0 & 0 \\
1 & 1 \\
2 & 2 \\
3 & 4 \\
4 & 10 \\
5 & 34 \\
6 & 154 \\
7 & 874 \\
8 & 5914 \\
9 & 46234 \\
10 & 409114 \\
\bottomrule
\end{tabular}
\end{center}

Zur Überprüfung rechnen wir einige Werte von Hand nach:
\begin{align*}
  !0 &= 0,\\
  !1 &= 0! = 1,\\
  !2 &= 0! + 1! = 1 + 1 = 2,\\
  !3 &= 0! + 1! + 2! = 1 + 1 + 2 = 4,\\
  !4 &= 0! + 1! + 2! + 3! = 1 + 1 + 2 + 6 = 10,\\
  !5 &= 1 + 1 + 2 + 6 + 24 = 34,\\
  !6 &= 34 + 5! = 34 + 120 = 154,\\
  !7 &= 154 + 6! = 154 + 720 = 874.
\end{align*}

\subsection{Rekursionsformel}

Die Links-Faktori\-elle genügt der einfachen Rekursion
\begin{equation}\label{eq:rekursion}
  !(n+1) \;=\; !n + n!\,, \qquad !0 = 0\,.
\end{equation}
Dies folgt unmittelbar aus der Definition:
$!(n+1) = \sum_{k=0}^{n} k! = \sum_{k=0}^{n-1} k! + n! = !n + n!$.

\subsection{Integraldarstellung}

Die Links-Faktori\-elle besitzt eine elegante Integraldarstellung.
Kurepa stellte fest, dass für $n \geq 1$ gilt:
\begin{equation}\label{eq:integral}
  !n \;=\; \int_0^{\infty} \frac{e^{-t}(t^n - 1)}{t - 1}\,dt\,,
\end{equation}
wobei der Integrand an der Stelle $t = 1$ durch den Grenzwert
(de l'H\^{o}pital) zu interpretieren ist.

\begin{proof}
Für $n \geq 1$ gilt:
\begin{align*}
  \int_0^{\infty} \frac{e^{-t}(t^n - 1)}{t - 1}\,dt
  &= \int_0^{\infty} e^{-t} \sum_{k=0}^{n-1} t^k \,dt
  = \sum_{k=0}^{n-1} \int_0^{\infty} t^k e^{-t}\,dt
  = \sum_{k=0}^{n-1} k! = !n\,. \qedhere
\end{align*}
Dabei haben wir die geometrische Summe
$(t^n - 1)/(t-1) = \sum_{k=0}^{n-1} t^k$ sowie die Laplace-Darstellung
$k! = \int_0^\infty t^k e^{-t}\,dt$ verwendet.
Der Austausch von Summe und Integral ist wegen der Dominanz durch
$n \cdot e^{-t/2}$ für große $t$ gerechtfertigt.
\end{proof}

\subsection{Wachstum}

Da der letzte Summand $(n-1)!$ alle anderen dominiert, gilt nach Stirling:
\[
  !n \;\sim\; (n-1)! \quad \text{für } n \to \infty\,.
\]
Genauer ist $!n = (n-1)!\bigl(1 + O(1/(n-1))\bigr)$.

\subsection{Teilbarkeit}

\begin{proposition}
Für $n \geq 3$ ist $!n$ gerade.
\end{proposition}
\begin{proof}
$!n = 1 + 1 + 2 + \sum_{k=3}^{n-1} k!$.
Für $k \geq 2$ gilt $k! \equiv 0 \pmod{2}$, also
$!n \equiv 1 + 1 + 0 = 2 \equiv 0 \pmod{2}$.
\end{proof}

\begin{bemerkung}
Für kleine Primzahlen $q$ zeigen die Reste $!n \bmod q$ kein offensichtliches
Muster; die Kurepa-Vermutung betrifft den Fall $q = p$, wobei $n = p$ selbst
eine Primzahl ist.
\end{bemerkung}

% ---------------------------------------------------------------
\section{Die Kurepa-Vermutung}
\label{sec:vermutung}
% ---------------------------------------------------------------

\subsection{Formulierung}

\begin{vermutung}[Kurepa, 1971]\label{verm:kurepa}
Für jede ungerade Primzahl $p$ gilt
\[
  p \nmid\; !p\,,
\]
d.\,h.\ $\displaystyle\sum_{k=0}^{p-1} k! \;\not\equiv\; 0 \pmod{p}$.
\end{vermutung}

\begin{bemerkung}
Die Einschränkung auf \emph{ungerade} Primzahlen ist wesentlich: für $p = 2$
gilt $!2 = 2 \equiv 0 \pmod{2}$, also ist $p = 2$ die einzige Primzahl,
für die die Behauptung scheitert.
\end{bemerkung}

\subsection{Der Ausnahmefall $p = 2$}

Für $p = 2$ ergibt sich:
\[
  !2 = 0! + 1! = 1 + 1 = 2 \equiv 0 \pmod{2}\,.
\]
Somit ist $p = 2$ tatsächlich ein Gegenbeispiel zur naiven Aussage ohne die
Einschränkung auf ungerade Primzahlen.
Die Hypothese ``ungerade Primzahl'' in Vermutung~\ref{verm:kurepa} ist daher
notwendig.

Für die ersten ungeraden Primzahlen:
\begin{align*}
  p = 3 &: \quad !3 = 4 \equiv 1 \pmod{3}\,, \quad 3 \nmid !3\,, \checkmark\\
  p = 5 &: \quad !5 = 34 \equiv 4 \pmod{5}\,, \quad 5 \nmid !5\,, \checkmark\\
  p = 7 &: \quad !7 = 874 = 124 \cdot 7 + 6 \equiv 6 \pmod{7}\,, \quad 7 \nmid !7\,. \checkmark
\end{align*}

\subsection{Kurepa-Primzahlen}

Falls eine ungerade Primzahl $p$ mit $p \mid !p$ existieren würde, nennt man
sie eine \emph{Kurepa-Primzahl}.
Die Kurepa-Vermutung behauptet, dass keine Kurepa-Primzahl existiert.
Bisher wurde trotz umfangreicher Computersuche noch nie eine gefunden.

% ---------------------------------------------------------------
\section{Verbindung zu Wilsons Satz}
\label{sec:wilson}
% ---------------------------------------------------------------

\subsection{Wilsons Satz}

Wir erinnern an den klassischen Satz:

\begin{theorem}[Wilson, 1770; Beweis durch Lagrange, 1771]\label{satz:wilson}
Eine ganze Zahl $n \geq 2$ ist genau dann prim, wenn
\[
  (n-1)! \;\equiv\; -1 \pmod{n}\,.
\]
\end{theorem}

Wilsons Satz verbindet die gewöhnliche Faktori\-elle mit der Primzahleigenschaft.
Die Kurepa-Vermutung schafft über die Links-Faktori\-elle eine zweite,
subtilere Verbindung.

\subsection{Reduktion von $!p \pmod{p}$}

Mit der Rekursion $!p = !(p-1) + (p-1)!$ und Wilsons Satz ergibt sich:

\begin{proposition}\label{prop:wilson-reduktion}
Für jede ungerade Primzahl $p$ gilt
\[
  !p \;\equiv\; !(p-1) - 1 \pmod{p}\,.
\]
\end{proposition}
\begin{proof}
Aus der Rekursion~\eqref{eq:rekursion} folgt
$!p = !(p-1) + (p-1)!$.
Nach Wilsons Satz (Satz~\ref{satz:wilson}) gilt
$(p-1)! \equiv -1 \pmod{p}$.
Daher $!p \equiv !(p-1) + (-1) = !(p-1) - 1 \pmod{p}$.
\end{proof}

\subsection{Äquivalente Formulierung über $!(p-1)$}

\begin{korollar}\label{kor:aequiv}
Die Kurepa-Vermutung ist äquivalent zur Aussage:
\[
  !(p-1) \;\not\equiv\; 1 \pmod{p}
\]
für jede ungerade Primzahl $p$.
\end{korollar}
\begin{proof}
Nach Proposition~\ref{prop:wilson-reduktion} gilt
$!p \equiv !(p-1) - 1 \pmod{p}$.
Daher gilt $p \mid !p$ genau dann, wenn $!(p-1) \equiv 1 \pmod{p}$.
\end{proof}

Diese Umformulierung ist praktisch: statt $p$ Summanden zu akkumulieren,
summiert man nur $p-1$ Summanden und prüft, ob das Ergebnis gerade $1$ modulo
$p$ ist.

\subsection{Stabile Reste}

\begin{proposition}
Sei $p$ eine ungerade Primzahl. Dann gilt für alle $m \geq p$:
\[
  !m \;\equiv\; !p \pmod{p}\,.
\]
\end{proposition}
\begin{proof}
Für $k \geq p$ enthält $k!$ den Faktor $p$, also $k! \equiv 0 \pmod p$.
Daher trägt jeder Term $k!$ mit $k \geq p$ nichts zur Summe $!m$ modulo $p$
bei.
\end{proof}

Der Rest $!p \bmod p$ ist also der eindeutige ``stabile Wert'' der Links-Fakto\-rial-Folge modulo $p$.

% ---------------------------------------------------------------
\section{Probabilistische Heuristik}
\label{sec:heuristik}
% ---------------------------------------------------------------

\subsection{Gleichverteilungshypothese}

Eine natürliche Annahme lautet: Die Reste $!p \bmod p$ verhalten sich
``zufällig'', d.\,h. sie sind näherungsweise gleichverteilt in
$\{0, 1, \ldots, p-1\}$, wenn $p$ über alle Primzahlen läuft.

Unter dieser Annahme beträgt die Wahrscheinlichkeit, dass eine gegebene
Primzahl $p$ eine Kurepa-Primzahl ist, ungefähr $1/p$.

\subsection{Erwartete Anzahl von Gegenbeispielen}

\begin{proposition}[Heuristik]
Unter der Gleichverteilungsannahme ist die erwartete Anzahl der
Kurepa-Primzahlen bis $x$:
\[
  \sum_{\substack{p \leq x \\ p \text{ ungerade Primzahl}}} \frac{1}{p}
  \;\sim\; \log\log x\,,
\]
was für $x \to \infty$ gegen unendlich divergiert.
\end{proposition}

Dies folgt aus dem Mertens'schen Satz
$\sum_{p \leq x} 1/p \sim \log\log x$.
Die Divergenz dieser Reihe bedeutet: \emph{heuristisch} sollten unendlich
viele Kurepa-Primzahlen existieren.

\subsection{Das Paradoxon}

Die heuristische Schlussfolgerung steht in scharfem Widerspruch zur
Kurepa-Vermutung und zur numerischen Evidenz:
Unter den ersten $664{.}578$ Primzahlen bis $10^7$ wurde keine einzige
Kurepa-Primzahl gefunden.

Ähnliche Paradoxa kennt man in der Zahlentheorie:
Beim Skewes-Problem erwartet man unendlich viele Kreuzungen von $\pi(x)$ und
$\mathrm{Li}(x)$, doch die erste liegt jenseits astronomisch großer Zahlen.
Bei der Collatz-Vermutung legen probabilistische Argumente nahe, dass alle
Folgen zu 1 konvergieren, was mit dem numerischen Befund übereinstimmt.

\subsection{Warum die Heuristik versagen könnte}

Die Gleichverteilungsannahme setzt implizit die Unabhängigkeit der Reste
$!p \bmod p$ voraus.
Die Werte $!p \bmod p$ werden jedoch durch eine algebraische Formel bestimmt,
die von der gesamten arithmetischen Struktur von $p$ abhängt.
Es könnte tiefe strukturelle Gründe geben --- in der Gruppentheorie, der
$p$-adischen Analysis oder der Arithmetik endlicher Körper ---,
die das Auftreten von $!p \equiv 0 \pmod p$ verhindern.

Insbesondere deutet die Verbindung zu Wilsons Satz (Abschnitt~\ref{sec:wilson})
darauf hin, dass $!p \bmod p$ arithmetische Information über $p$ kodiert,
die ein rein probabilistisches Modell nicht erfasst.

% ---------------------------------------------------------------
\section{Äquivalenzen und Umformulierungen}
\label{sec:aequivalenzen}
% ---------------------------------------------------------------

\subsection{Zusammenfassung der Äquivalenzen}

Die folgenden Aussagen sind äquivalent:
\begin{enumerate}[label=(\alph*)]
\item (Kurepa) $p \nmid !p$ für jede ungerade Primzahl $p$.
\item $!(p-1) \not\equiv 1 \pmod{p}$ für jede ungerade Primzahl $p$.
\item Für jede ungerade Primzahl $p$ existiert ein $r \in \{1,\ldots,p-1\}$
      mit $\sum_{k=0}^{p-1} k! \equiv r \pmod{p}$.
\end{enumerate}

Die Äquivalenz von (a) und (b) ist Korollar~\ref{kor:aequiv};
(c) ist eine Umformulierung von (a).

\subsection{Verbindung zur Subfaktoriellen}

Die Subfaktorielle $D_n = n! \sum_{k=0}^n (-1)^k/k!$ zählt die
Derangements von $\{1,\ldots,n\}$.
Zwischen $!n$ und $D_n$ besteht keine einfache geschlossene Beziehung im
Allgemeinen; beide Folgen sind jedoch $D$-endlich (d.\,h.\ sie genügen
linearen Rekursionen mit Polynomkoeffizienten).

\subsection{Gruppentheoretischer Zugang}

Sei $G = (\Z/p\Z)^*$ die multiplikative Gruppe modulo $p$ der Ordnung $p-1$.
Die Summe $!p = \sum_{k=0}^{p-1} k! \pmod p$ kann im Gruppenring $\Z[G]$
durch Umschreiben von $k!$ mit Hilfe von Wilson-Quotienten ausgedrückt werden.
Diese algebraische Perspektive könnte gruppentheoretische Methoden erschliessen.

\subsection{Verbindung zum Brocard--Ramanujan-Problem}

Das Brocard--Ramanujan-Problem fragt, für welche $n$ die Gleichung
$n! + 1 = m^2$ eine Lösung $m \in \N$ hat.
Bekannte Lösungen: $n \in \{4, 5, 7\}$.
Beide Probleme betreffen Faktori\-elle unter arithmetischen
Divisibilitätsbedingungen, ohne dass eine direkte Reduktion zwischen ihnen
bekannt wäre.

\subsection{Bernoullizahlen}

Die Summe $\sum_{k=0}^{p-1} k!$ lässt sich formal in Verbindung mit
faktorgewichteten Partialsummen bringen.
Eine explizite Darstellung über Bernoullizahlen ist nicht bekannt, doch
Kurepas Originalpaper erwähnt Verbindungen zu kombinatorischen Identitäten
modulo Primzahlen.

% ---------------------------------------------------------------
\section{Numerische Verifikation}
\label{sec:verifikation}
% ---------------------------------------------------------------

\subsection{Beschreibung der Berechnung}

Das Specialist Mathematics Project (Build 131) hat die
Kurepa-Vermutung~\ref{verm:kurepa} systematisch numerisch verifiziert.
Der Algorithmus besteht aus folgenden Schritten:

\begin{enumerate}
\item Erzeuge alle Primzahlen $p \leq N$ mit dem Sieb des Eratosthenes.
\item Berechne für jede ungerade Primzahl $p$ den Wert $!p \bmod p$
      mittels der Rekursion
      $!(k+1) \equiv !k + k! \pmod p$,
      wobei $k!$ und $!k$ als laufende Werte modulo $p$ geführt werden.
\item Notiere $p$ als Verletzung, falls $!p \equiv 0 \pmod p$.
\end{enumerate}

In Schritt~2 muss $k!$ niemals exakt berechnet werden; stattdessen werden
der laufende Faktor $k! \bmod p$ und die laufende Summe $!k \bmod p$
gemeinsam aktualisiert.
Dies ergibt einen effizienten $O(p)$-Algorithmus für jede Primzahl $p$.

\subsection{Ergebnisse}

\begin{center}
\begin{tabular}{@{}ll@{}}
\toprule
Parameter & Wert \\
\midrule
Obere Schranke $N$ & $10^7$ \\
Anzahl getesteter Primzahlen & $664{.}578$ \\
Anzahl gefundener Verletzungen & $0$ \\
Rechenzeit & $\approx 8{.}396$ Sekunden \\
Build & 131 (Specialist Mathematics Project) \\
\bottomrule
\end{tabular}
\end{center}

\begin{bemerkung}
Die Zahl $664{.}578$ ist die Anzahl der ungeraden Primzahlen bis $10^7$;
die Gesamtzahl der Primzahlen bis $10^7$ beträgt $\pi(10^7) = 664{.}579$,
wobei $p = 2$ herausgenommen wird.
Das Ergebnis lautet: \emph{Es gibt keine Kurepa-Primzahl unterhalb von $10^7$}.
\end{bemerkung}

\subsection{Historische Verifikationen im Vergleich}

\begin{center}
\begin{tabular}{@{}lll@{}}
\toprule
Prüfer & Schranke & Jahr \\
\midrule
Kurepa & Kleine Fälle & 1971 \\
Stankov & $p \leq 10^6$ & 2003 \\
Verschiedene & $p \leq 10^8$ & nach 2010 \\
Specialist Mathematics Project (Build 131) & $p \leq 10^7$ & 2026 \\
\bottomrule
\end{tabular}
\end{center}

\subsection{Implementierungshinweise}

Die Python-Implementierung verwendet durchgehend modulare Arithmetik,
um den Umgang mit großen Ganzzahlen zu vermeiden.
Entscheidend ist, dass für $k \geq p$ gilt $k! \equiv 0 \pmod p$,
sodass die Schleife effektiv nur $p$ Iterationen durchläuft.
Der Engpass liegt beim Sieb des Eratosthenes für $N = 10^7$ sowie in der
anschließenden modularen Akkumulationsschleife.

\begin{beispiel}
Für $p = 7$:
\begin{align*}
  !7 &= 0! + 1! + 2! + 3! + 4! + 5! + 6!\\
     &= 1 + 1 + 2 + 6 + 24 + 120 + 720 = 874\,.
\end{align*}
Modulo 7: $874 = 124 \cdot 7 + 6$, also $!7 \equiv 6 \pmod 7 \neq 0$.
Die Vermutung gilt für $p = 7$.
\end{beispiel}

\begin{beispiel}
Für $p = 11$:
\[
  !11 = 0! + 1! + \cdots + 10!\,.
\]
Nach Wilsons Satz gilt $10! \equiv -1 \pmod{11}$.
Eine Rechnung ergibt $!11 \equiv 1 \pmod{11} \neq 0$.
\end{beispiel}

\begin{beispiel}
Für $p = 13$:
\[
  !13 = \sum_{k=0}^{12} k!\,.
\]
Eine Berechnung modulo 13 liefert $!13 \equiv 9 \pmod{13} \neq 0$.
\end{beispiel}

% ---------------------------------------------------------------
\section{Verwandte Probleme}
\label{sec:verwandt}
% ---------------------------------------------------------------

\subsection{Wilson-Primzahlen}

Der Wilson-Quotient einer Primzahl $p$ ist $W_p = ((p-1)! + 1)/p$.
Eine \emph{Wilson-Primzahl} ist eine Primzahl $p$ mit $p \mid W_p$,
gleichbedeutend mit $(p-1)! \equiv -1 \pmod{p^2}$.
Bisher sind nur drei Wilson-Primzahlen bekannt: $5, 13, 563$.
Die Kurepa-Vermutung und die Theorie der Wilson-Primzahlen sind durch die
gemeinsame Nutzung von $(p-1)! \bmod p$ verbunden, ohne dass eine direkte
Reduktion zwischen beiden Problemen bekannt wäre.

\subsection{$p$-adische Perspektive}

Die $p$-adische Bewertung $v_p(!p)$ beträgt genau dann $0$,
wenn die Kurepa-Vermutung für $p$ gilt (da $!p$ dann eine $p$-adische Einheit
ist).
Man könnte versuchen, durch $p$-adische Methoden zu zeigen, dass $!p$ stets
ausserhalb des maximalen Ideals von $\Z_p$ liegt.

\subsection{Verallgemeinerungen der Links-Faktoriellen}

Man kann die Links-Faktori\-elle verallgemeinern:
\[
  !_m n \;=\; \sum_{k=0}^{n-1} (mk)!\,.
\]
Das Verhalten solcher Verallgemeinerungen modulo Primzahlen ist weitgehend
unerforscht.

\subsection{Das Brocard--Ramanujan-Problem}

Dieses Problem fragt nach ganzzahligen Lösungen von $n! + 1 = m^2$.
Bisher bekannte Lösungen: $(n, m) = (4, 5), (5, 11), (7, 71)$.
Ob weitere Lösungen existieren, ist offen.
Ähnlich wie die Kurepa-Vermutung verknüpft dieses Problem
Faktorielle mit Divisibilitätsbedingungen, ohne dass eine Reduktion zwischen
beiden bekannt wäre.

\subsection{Faktorielle Primzahlen}

Eine \emph{faktorielle Primzahl} ist eine Primzahl der Form $n! \pm 1$.
Diese sind ebenfalls mit Partialsummen von Faktorielle verwandt,
stehen aber in keiner bekannten direkten Beziehung zur Kurepa-Vermutung.

\subsection{Smarandache-Funktion}

Die Smarandache-Funktion $S(n)$ ist das kleinste $m$, sodass $n \mid m!$.
Sowohl $S(n)$ als auch $!n$ verknüpfen Faktorielle mit Teilbarkeitsfragen,
ohne dass eine bekannte direkte Implikation zwischen der Kurepa-Vermutung und
Aussagen über $S$ besteht.

% ---------------------------------------------------------------
\section{Schluss}
\label{sec:schluss}
% ---------------------------------------------------------------

Die Kurepa-Vermutung ist eines der elementarsten offenen Probleme der
Zahlentheorie.
Ihre Formulierung erfordert nichts weiter als die Definition der Faktoriellen
und modulare Arithmetik, und dennoch hat sie seit 1971 allen Beweisversuchen
widerstanden.

Die wesentlichen Aspekte des Problems lassen sich so zusammenfassen:

\begin{enumerate}
\item \textbf{Elementarität}: Die Aussage ist für jeden Studierenden verständlich,
      der die modulare Arithmetik kennt.
\item \textbf{Wilson-Verbindung}: Die Formel $!p \equiv !(p-1) - 1 \pmod p$
      verbindet die Vermutung mit Wilsons klassischem Satz und reduziert sie
      auf $!(p-1) \not\equiv 1 \pmod p$.
\item \textbf{Heuristisches Paradoxon}: Probabilistische Argumente sagen
      unendlich viele Gegenbeispiele voraus, doch keines wurde je gefunden.
\item \textbf{Numerische Evidenz}: Null Verletzungen unter $664{.}578$
      ungeraden Primzahlen bis $10^7$ (Build 131).
\end{enumerate}

\subsection*{Offene Fragen}
\begin{enumerate}
\item Gibt es eine ungerade Primzahl $p$ mit $p \mid !p$? (Kurepa-Vermutung)
\item Ist $!p \bmod p$ gleichverteilt in $\{1, \ldots, p-1\}$ für $p \to \infty$?
\item Gibt es einen $p$-adischen oder algebraischen Beweisansatz,
      der ohne Computerrechnung auskommt?
\item Was ist die genaue Größenordnung von $\min_{p \text{ prim}} v_p(!p)$
      als Funktion von $p$?
\end{enumerate}

Die in Build 131 durchgeführte numerische Verifikation erweitert die bekannte
Schranke für Nicht-Gegenbeispiele und bildet eine konkrete Grundlage für
weitere theoretische Untersuchungen.
Wir hoffen, dass diese Darstellung das Problem einem breiteren Publikum
zugänglich macht und neue Beweisideen anregt.

% ---------------------------------------------------------------
% Literaturverzeichnis
% ---------------------------------------------------------------
\begin{thebibliography}{99}

\bibitem{Kurepa1971}
\DJ{}. Kurepa,
\textit{On the left factorial function $!n$},
Math.\ Balkanica \textbf{1} (1971), 147--153.

\bibitem{Stankov2003}
D. Stankov,
\textit{Über das Kurepa-Problem},
Filomat \textbf{17} (2003), 61--65.

\bibitem{Wilson1770}
J. Wilson (mitgeteilt durch E. Waring),
\textit{Brief an E. Waring},
1770; veröffentlicht in E. Waring, \textit{Meditationes Algebraicae},
Cambridge, 1770.

\bibitem{Lagrange1771}
J.-L. Lagrange,
\textit{Demonstration d'un th\'eor\`eme nouveau concernant les nombres premiers},
Nouveaux M\'emoires de l'Acad\'emie Royale des Sciences et Belles-Lettres de Berlin,
1771.

\bibitem{Mertens1874}
F. Mertens,
\textit{Ein Beitrag zur analytischen Zahlentheorie},
J.\ reine angew.\ Math.\ \textbf{78} (1874), 46--62.

\bibitem{GuyUnsolved}
R.\ K.\ Guy,
\textit{Unsolved Problems in Number Theory}, 3.\ Aufl.,
Springer, New York, 2004.
Problem B44.

\bibitem{SloaneA002627}
N.\ J.\ A.\ Sloane u.\,a.,
\textit{The On-Line Encyclopedia of Integer Sequences},
Folge A002627 (Links-Faktori\-elle $!n$),
\url{https://oeis.org/A002627}.

\bibitem{SpecialistBuild131}
M.\ Fuhrmann,
\textit{Specialist Mathematics Project, Build 131: Numerische Verifikation der
Kurepa-Vermutung für Primzahlen bis $10^7$},
Interner Bericht, 2026.

\end{thebibliography}

\end{document}
